\section{Introduction}

Music understanding systems increasingly combine several kinds of input: a natural-language
description of a track, and the track's own audio. Audio is often turned into a graph before
it reaches a model -- a track is split into short segments, each becomes a node, and edges
connect segments close in time or acoustically similar. This is a reasonable design choice:
music has real relational structure (repetition, transition, similarity between sections),
and a graph neural network (GNN) can in principle use it. Similarity between short audio
segments, computed from acoustic features and thresholded or visualized as a matrix, is a
long-established technique for finding structure within a track \cite{foote2000novelty};
this paper's graph construction applies the same idea (pairwise segment similarity) to build
a graph rather than a novelty curve. The message-passing GNN family this paper draws its
encoder from includes several related architectures beyond the one used here -- GCN
\cite{kipf2017gcn}, GAT \cite{velickovic2018gat}, and GIN \cite{xu2019gin}, alongside
GraphSAGE \cite{hamilton2017graphsage}. Separately, a recognized concern in the broader GNN
literature is that graph-based models are not always compared against simpler baselines
under matched, controlled conditions \cite{errica2020fair,shchur2018pitfalls} -- this paper
takes that concern as a direct methodological starting point, not only a citation.

But if the edges are built directly from the same features already given to each node --
connecting two segments whenever their feature similarity crosses a threshold -- then the
graph's structure is, information-theoretically, a restatement of information already
present in the node features, not a new observation about the track. This does not
automatically mean such a graph is useless: a deterministic function of the features could
still act as a useful computational shortcut or inductive bias, easing optimization even
without adding information a sufficiently powerful model could not already extract from the
features alone. Whether either possibility -- added information, or a
useful-but-informationless shortcut -- actually helps in practice, for a specific
construction and specific tasks, is an empirical question, not something to assume from the
architecture diagram alone.

This paper answers that question for one concrete system, across four tasks sharing the same
node-feature definition and similarity-based edge construction rule -- though not, as
Section~\ref{sec:data} discloses precisely, the same audio segmentation length or dataset
split in every task: a text-only tag classifier (Task 1, no audio at all), a graph-based
genre classifier compared against an edge-free alternative (Task 2), a graph-plus-text fusion
model (Task 3), and a contrastive audio-text retrieval model (Task 4). Each task, alone,
answers a narrower question. Together, they form a single experimental ladder:

\begin{itemize}
\item \textbf{Task 1} asks how much target information language alone already provides,
before any audio or graph is involved.
\item \textbf{Task 2} asks whether the graph's topology helps audio-only classification
beyond what the same node features give a structure-free model.
\item \textbf{Task 3} asks whether the graph adds information that a strong text model does
not already have, once both are available together.
\item \textbf{Task 4} asks whether any graph benefit found in classification or fusion
survives under a different objective entirely: matching audio to text by similarity, with no
genre or tag labels -- only audio-caption pairing as supervision.
\end{itemize}

We do not set out to show that graphs are useless for music, and we do not find that. We set
out to check, directly and with matched controls wherever they could be arranged --
Section~\ref{sec:protocol} discloses the specific exceptions found on inspection, including
one confound in Task 2's own training scripts -- whether \emph{this} similarity-based graph
construction earns its place in \emph{these} tasks. Task 1 has no graph and serves as the
anchor; Tasks 2--4 each include a direct graph-vs-edge-free comparison, and in none does the
graph-based model come out clearly and reliably ahead. Task 2's fairest (pooling-matched)
comparison favors the edge-free alternative, though a non-learned 2-hop control complicates
this (Section~\ref{sec:discussion}). Task 3's standalone comparison is mixed and
inconclusive, but its two follow-up diagnostics tell distinct, narrower stories:
branch-zeroing (an inference-time intervention on the trained fusion model) finds it relies
only slightly on the graph branch -- topology and node features together, not
topology-isolating -- while complementarity (correlational, comparing two separately-trained
models rather than probing the fusion model) finds almost no positive-label rescue by the
graph-only model, no sign it detects tags text misses. Together, in this evidence base, they
point toward little demonstrated graph value for Task 3, though neither is topology-isolating
and complementarity says nothing about the fusion model's own behavior; a follow-up check in
a second implementation did not confirm even this narrower reading, reported openly in
Section~\ref{sec:discussion} rather than alongside only the reading that fits. Task 4 shows
no consistent advantage either way across three seeds. We report this consistency, with each
result's own caveats, as the paper's main finding. We also check two alternative explanations
that could produce this pattern without the graph itself being at fault -- caption text
directly leaking target labels, and duplicate or leaked audio inflating a model's score --
and report what we find, rather than assume either away.
